\documentclass[12pt,a4paper]{article}
\usepackage[margin=2cm]{geometry}
\usepackage{graphicx}
\usepackage{tabularx}
\usepackage{array}
\usepackage{enumitem}
\usepackage{hyperref}
\usepackage{float}

% Unicode friendly (XeLaTeX)
\usepackage{fontspec}
\defaultfontfeatures{Ligatures=TeX}

% For checkbox symbol
\usepackage{amssymb}
\providecommand{\checkbox}{$\square$}

% For grids
\usepackage{tikz}
\usetikzlibrary{calc}

\setlength{\parindent}{0pt}
\setlength{\parskip}{6pt}

% A line field
\newcommand{\linefield}[1]{\rule{#1}{0.4pt}}

% Task heading
\newcommand{\tasktitle}[2]{\vspace{6pt}\textbf{#1}\hfill /#2\par\vspace{4pt}}

\begin{document}

\def\SchoolLogo{C:/Users/TKast/OneDrive/Dokumente/GitHub/SchulaufgabenScript/logos/krs.png}
\def\ExamTitle{1. Schulaufgabe}
\def\ExamSubject{Mathematik}
\def\ExamDate{2024-11-26}
\def\ExamClass{9}
\def\TotalPoints{12.5}

% ==========================================
% Schulaufgaben Header Template
% (external file loaded by Python script)
%
% Expected variables:
% \SchoolLogo
% \ExamTitle
% \ExamSubject
% \ExamDate
% \ExamClass
% \TotalPoints
% ==========================================

\newcommand{\cb}{$\square$}

\noindent
\setlength{\fboxsep}{0pt}
\setlength{\fboxrule}{0.9pt}

\fbox{%
\begin{minipage}{\linewidth}
\renewcommand{\arraystretch}{1.15}

\begin{tabular}{@{}p{0.20\linewidth}@{}|@{}p{0.80\linewidth}@{}}

% =================================================
% LEFT SIDE (LOGO + CHECKBOXES)
% =================================================
\begin{minipage}[t]{\linewidth}
\vspace{6pt}
\centering
\includegraphics[width=0.85\linewidth]{\SchoolLogo}

\vspace{10pt}
\hrule
\vspace{10pt}

\raggedright
\hspace{10pt}\cb\ LRSt\par\vspace{6pt}
\hspace{10pt}\cb\ IRSt\par\vspace{6pt}
\hspace{10pt}\cb\ ILSt\par\vspace{10pt}

\hspace{10pt}\cb\ \rule{2.3cm}{0.4pt}\par
\vspace{8pt}

\end{minipage}
&
% =================================================
% RIGHT SIDE
% =================================================
\begin{minipage}[t]{\linewidth}

% ---------- TITLE ROW ----------
\begin{tabularx}{\linewidth}{@{}|X|p{0.22\linewidth}|p{0.18\linewidth}|@{}}
\hline
\textbf{\ExamTitle} &
\textbf{Fach}\par \ExamSubject &
\textbf{Datum}\par \ExamDate \\
\hline
\end{tabularx}

% ---------- NAME / CLASS / NOTE ----------
\begin{tabularx}{\linewidth}{@{}|X|p{0.16\linewidth}|p{0.16\linewidth}|p{0.16\linewidth}|@{}}
\hline

\parbox[t][3.4cm][t]{\linewidth}{
\textbf{Name, Vorname}

\vfill
\rule{\linewidth}{0.4pt}
}
&
\parbox[t][3.4cm][t]{\linewidth}{
\textbf{Klasse}

\vfill
\rule{\linewidth}{0.4pt}
}
&
\parbox[t][3.4cm][t]{\linewidth}{
\textbf{Nr.}

\vfill
\rule{\linewidth}{0.4pt}
}
&
\parbox[t][3.4cm][t]{\linewidth}{
\textbf{Note}

\vfill
\rule{\linewidth}{0.4pt}
}

\\
\hline
\end{tabularx}

% ---------- SIGNATURE ----------
\begin{tabularx}{\linewidth}{@{}|X|p{0.20\linewidth}|@{}}
\hline

\parbox[t][3.4cm][b]{\linewidth}{
\textit{Unterschrift des/der Erziehungsberechtigten}

\vspace{6pt}
\rule{\linewidth}{0.4pt}
}
&
\\
\hline
\end{tabularx}

% ---------- TIME EXTENSION + POINTS ----------
\begin{tabularx}{\linewidth}{@{}|X|p{0.20\linewidth}|@{}}
\hline

\parbox[t][2.6cm][t]{\linewidth}{
Zeitzuschlag in Anspruch genommen:\par\vspace{8pt}

\cb\ ja \hspace{14pt}
\cb\ nein \hspace{14pt}
\cb\ teilweise

\vspace{10pt}
\rule{2.2cm}{0.4pt}\ Minuten
}
&
\parbox[t][2.6cm][c]{\linewidth}{
\centering
\Large \underline{\hspace{1.5cm}}/\TotalPoints
}

\\
\hline
\end{tabularx}

\end{minipage}
\\
\end{tabular}

\end{minipage}
}

\vspace{12pt}

\tasktitle{Aufgabe 1: Fehlerbilder zu linearen Funktionen}{4.5}
Verbinde jede Zeichnung mit der dazu passenden Beschreibung des Fehlers. Fehlt eine passende Beschreibung, formuliere sie im letzten Kasten.

\begin{center}\includegraphics[width=\linewidth]{C:/Users/TKast/OneDrive/Dokumente/GitHub/SchulaufgabenScript/tasks/lineare_funktionen_fehlerbilder_01/images/layout.png}\end{center}


\tasktitle{Aufgabe 2: Gerade g₁: Punktprobe, fehlende Koordinate, Nullstelle, Parallelität}{8.0}
Gegeben ist die Gerade $g_1: y = -\frac{2}{3}x + 2$.

\begin{enumerate}[label=\alph*)]
\item Prüfe rechnerisch, ob der Punkt $P\,(-1{,}5\mid 3)$ auf der Geraden $g_1$ liegt.
\begin{center}
\begin{tikzpicture}
  \draw[step=0.5cm, very thin] (0,0) grid (16.0cm, 4.0cm);
  \draw[line width=0.4pt] (0,0) rectangle (16.0cm, 4.0cm);
\end{tikzpicture}
\end{center}

\item Der Punkt $B\,(x\mid \frac{1}{3})$ liegt auf der Geraden $g_1$. Berechne die fehlende Koordinate.
\begin{center}
\begin{tikzpicture}
  \draw[step=0.5cm, very thin] (0,0) grid (16.0cm, 4.0cm);
  \draw[line width=0.4pt] (0,0) rectangle (16.0cm, 4.0cm);
\end{tikzpicture}
\end{center}

\item Berechne die Nullstelle der Geraden $g_1$.
\begin{center}
\begin{tikzpicture}
  \draw[step=0.5cm, very thin] (0,0) grid (16.0cm, 4.0cm);
  \draw[line width=0.4pt] (0,0) rectangle (16.0cm, 4.0cm);
\end{tikzpicture}
\end{center}

\item Prüfe rechnerisch, ob die Gerade $g_2$ mit der Gleichung $x + \frac{1}{3}y - 2 = 0$ parallel zur Geraden $g_1$ verläuft.
\begin{center}
\begin{tikzpicture}
  \draw[step=0.5cm, very thin] (0,0) grid (16.0cm, 5.0cm);
  \draw[line width=0.4pt] (0,0) rectangle (16.0cm, 5.0cm);
\end{tikzpicture}
\end{center}

\end{enumerate}
\vspace{6pt}\hrule\vspace{6pt}
\end{document}